%% Copyright (C) 2010, 2011, Gostai S.A.S.
%%
%% This software is provided "as is" without warranty of any kind,
%% either expressed or implied, including but not limited to the
%% implied warranties of fitness for a particular purpose.
%%
%% See the LICENSE file for more information.

\documentclass[openright,twoside,11pt]{book}
%  \usepackage[svgnames]{xcolor}  \usepackage{experiment}
  \usepackage{gostai-documentation}
  \usepackage{config}
  \usepackage[draft]{showlabels}
  \usepackage{lipsum}  % to provide "filler" text

\title{Urbi Documentation Sandbox}
\subtitle{Version \VcsDescription}
\author{Gostai}

\begin{document}

\maketitle

This document is meant to be used to toy with \LaTeX{}.

\tableofcontents
\part{First Part}

\chapter{First Chapter}

This is the first chapter.

\begin{verbatim}
for (var i: 10)
  echo(i);
\end{verbatim}

\begin{urbiunchecked}
'if' + 'else' + 'bitand';
'if'+'else'+'bitand';
'if'.'else'.'bitand';
.'else';

x xor_eq a;

a.'bitand'(b);
a' "foo";
a'' "foo";
a''' "foo";
"else";

"foo";
[00048238] "foo"
\end{urbiunchecked}

\begin{verbatim}
for (var i: 10)
  echo(i);
\end{verbatim}

\begin{verbatim}[language=gdb]
(gdb) urbi break input.u:2 Lobby_0x7ffff7f03208.foo(["2" => 2])
UBreakpoint 1:
        Location: input.u:2
        Lobby_0x7ffff7f03208.foo(["2" => 2])
(gdb) c
Continuing.

//#push 1 "input.u"
function foo (x) { backtrace }|;
for(var i: [1, 2, 3]) foo([i.asString => i]);

[00219506:backtrace] foo (input.u:2.23-44)
[00219506:backtrace] each (input.u:2.12-44)
UBreakpoint 1: [input.u:2.23-45] Lobby_0x7ffff7f03208.foo(["2" => 2])
Inside Job 0x6a2f70
(gdb) urbi stack
#12 [input.u:1.20-29] Lobby_0x7ffff7f03208.backtrace()
#16 [input.u:1.20-29] Call backtrace
#55 [input.u:2.1-40] Stmt
(gdb) #% comment
\end{verbatim}

\begin{bnf}
<class-statement>
  ::= "class" <lvalue> (: <prototypes>)? <block>

<lvalue>
  ::= (<expression> ".")* "identifier"

<foo>
  ::= "(" "*" ")"

<prototypes>>
  ::= (<expression> ",")* expression

<block>
  ::= "{" <statement>* "}"
\end{bnf}

Another example.

\begin{shell}[style=varInString]
$ export PATH="something:\var{urbi-root}/bin:$PATH"
$ test/then
\end{shell}%$

\begin{urbiunchecked}[style=varInString]
"for \var{foo}" * (1+2);
[00048238] wrong
\end{urbiunchecked}

Last example.

\begin{urbiassert}
1+2 == 3;
\end{urbiassert}

% \release{Release}{Date}
% -----------------------
\newcommand{\release}[2]{%
  \ifx\ifHtml\undefined%
    Released on #2.
    \marginpar{%
      \includegraphics[width=\marginparwidth]{figs/urbi-sdk/#1/cube}

      #2.
    }%
  \else
    \includegraphics[width=.2\textwidth]{figs/urbi-sdk/#1/cube},
    released on #2.
  \fi%
}

\section{\usdk 2.5}
\release{2.5}{2010-11-xx}

\subsection{New Features}
\begin{itemize}
\item More types of empty statements are warned about.  For instance \urbi
  used to accept silently \lstinline|if (foo);|.  It now warns about the empty
  body, and recommends writing \lstinline|if (foo) {};|.
\end{itemize}

\section{\usdk 2.4}
\release{2.4}{2010-11-xx}

\subsection{New Features}
\begin{itemize}
\item More types of empty statements are warned about.  For instance \urbi
  used to accept silently \lstinline|if (foo);|.  It now warns about the empty
  body, and recommends writing \lstinline|if (foo) {};|.
\end{itemize}

\section{\usdk 2.3}
\release{2.3}{2010-11-xx}

\subsection{New Features}
\begin{itemize}
\item More types of empty statements are warned about.  For instance \urbi
  used to accept silently \lstinline|if (foo);|.  It now warns about the empty
  body, and recommends writing \lstinline|if (foo) {};|.
\end{itemize}


\sectionObject{Bar}

\begin{urbiscriptapi}
\item[bar] This is Bar.bar.
\item[foo](<x>) This is Bar.foo with arguments.
\item[set](<key>, <value>) Map \var{key} to \var{value} and return \this so
  that invocations to \lstinline|set| can be chained.  The possibly existing
  previous mapping is overridden.
\item \lstinline|I do it myself|
\item beware (not to catch the first word)
\item|slot[foo][bar]| a slot
\item|fun[foo][bar]|(<var>) a function
\end{urbiscriptapi}

\sectionObject{Foo}

\begin{urbiscriptapi}
\item[bar] This is Foo.bar.
\item[foo] This is Foo.foo.
\end{urbiscriptapi}

\section{Second Section}

\begin{urbiunchecked}[escapeinside=<>]
#<Été>;
[00048238:error] !!! invalid character: `#'
[00048239:error] !!! invalid character: `\xc3'
[00048239:error] !!! invalid character: `\x89'
[00048239:error] !!! invalid character: `\xc3'
[00048239:error] !!! invalid character: `\xa9'
\end{urbiunchecked}

\chapter{Second Chapter}

This is the second chapter.

\urbitrajectory{smooth}

\urbitrajectory{accel}

\begin{urbiscriptapi}
\item[count] Return the count.
\item[launch]
  Fire \this.
\item \lstinline|launch|~\\
  Fire \this.
\end{urbiscriptapi}


\renewcommand{\baselinestretch}{.85}
\begin{table}[\floatpos]
  \centering
    \ifx\ifHtml\undefined
\begin{tabular}{|c|l||c|l|}
  \hline
  Keyword & Remark & Keyword & Remark \\
  \hline
     \lstinline"and" & Synonym for \lstinline|&&|               &
  \lstinline"and_eq" & Synonym for \lstinline|&=|               \\
     \lstinline"asm" & Reserved                                 &
  \lstinline"assert" & \autoref{sec:lang:assert}                \\
      \lstinline"at" & \autoref{sec:lang:at}                    &
    \lstinline"auto" & Reserved                                 \\
  \lstinline"bitand" & \autoref{sec:lang:op:bitwise}            &
   \lstinline"bitor" & \autoref{sec:lang:op:bitwise}            \\
    \lstinline"bool" & Reserved                                 &
   \lstinline"break" & \autoref{sec:lang:break}                 \\
    \lstinline"call" & \autoref{sec:lang:rv}                    &
    \lstinline"case" & \autoref{sec:lang:switch}                \\
   \lstinline"catch" & \autoref{sec:lang:catch}, \autoref{sec:lang:timeout} &
    \lstinline"char" & Reserved                                 \\
   \lstinline"class" & \autoref{sec:lang:class}                 &
 \lstinline"closure" & \autoref{sec:lang:closure}               \\
   \lstinline"compl" & \autoref{sec:lang:op:bitwise}            &
   \lstinline"const" & \autoref{sec:lang:prop:std}              \\
\lstinline"const_cast" & Reserved                                 &
\lstinline"continue" & \autoref{sec:lang:continue}              \\
 \lstinline"default" & \autoref{sec:lang:switch}                &
  \lstinline"delete" & Reserved                                 \\
  \lstinline"detach" &                                          &
  \lstinline"disown" &                                          \\
      \lstinline"do" & \autoref{sec:lang:do}                    &
  \lstinline"double" & Reserved                                 \\
\lstinline"dynamic_cast" & Reserved                                 &
    \lstinline"else" & \autoref{sec:lang:if}                    \\
    \lstinline"enum" & \autoref{sec:lang:continue}              &
   \lstinline"every" & \autoref{sec:lang:every}                 \\
\lstinline"explicit" & Reserved                                 &
  \lstinline"export" & Reserved                                 \\
  \lstinline"extern" & Reserved                                 &
\lstinline"external" & Reserved                                 \\
   \lstinline"float" & Reserved                                 &
     \lstinline"for" & \lstinline|for&| and \lstinline-for|- flavors \\
 \lstinline"foreach" & Deprecated, use \lstinline|for|          &
\lstinline"freezeif" &                                          \\
  \lstinline"friend" & Reserved                                 &
\lstinline"function" & \autoref{sec:lang:function}, \autoref{sec:lang:closure} \\
    \lstinline"goto" & Reserved                                 &
      \lstinline"if" & \autoref{sec:lang:if}                    \\
      \lstinline"in" & \autoref{sec:lang:foreach}               &
  \lstinline"inline" & Reserved                                 \\
     \lstinline"int" & Reserved                                 &
\lstinline"internal" & Deprecated                               \\
    \lstinline"long" & Reserved                                 &
    \lstinline"loop" & \lstinline|loop&|, \lstinline-loop|-, \autoref{sec:lang:loop} \\
 \lstinline"mutable" & Reserved                                 &
\lstinline"namespace" & Reserved                                 \\
     \lstinline"new" & \autoref{sec:tut:ctor}                   &
     \lstinline"not" & Synonym for \lstinline|!|                \\
  \lstinline"not_eq" & Synonym for \lstinline|!=|               &
 \lstinline"onleave" & \autoref{sec:lang:at}                    \\
      \lstinline"or" & Synonym for \lstinline-||-               &
   \lstinline"or_eq" & Synonym for \lstinline-|=-               \\
 \lstinline"private" & Ignored                                  &
\lstinline"protected" & Ignored                                  \\
  \lstinline"public" & Ignored                                  &
\lstinline"register" & Reserved                                 \\
\lstinline"reinterpret_cast" & Reserved                                 &
  \lstinline"return" & \autoref{sec:lang:rv}                    \\
   \lstinline"short" & Reserved                                 &
  \lstinline"signed" & Reserved                                 \\
  \lstinline"sizeof" & Reserved                                 &
  \lstinline"static" & Deprecated                               \\
\lstinline"static_cast" & Reserved                                 &
  \lstinline"stopif" &                                          \\
  \lstinline"struct" & Reserved                                 &
  \lstinline"switch" & \autoref{sec:lang:switch}                \\
\lstinline"template" & Reserved                                 &
    \lstinline"this" &                                          \\
   \lstinline"throw" & \autoref{sec:lang:throw}                 &
 \lstinline"timeout" & \autoref{sec:lang:timeout}               \\
     \lstinline"try" & \autoref{sec:lang:catch}                 &
 \lstinline"typedef" & Reserved                                 \\
  \lstinline"typeid" & Reserved                                 &
\lstinline"typename" & Reserved                                 \\
   \lstinline"union" & Reserved                                 &
\lstinline"unsigned" & Reserved                                 \\
   \lstinline"using" & Reserved                                 &
     \lstinline"var" & \autoref{sec:lang:var}, \autoref{sec:lang:slots} \\
 \lstinline"virtual" & Reserved                                 &
\lstinline"volatile" & Reserved                                 \\
\lstinline"waituntil" & \autoref{sec:lang:waituntil}             &
   \lstinline"watch" & \autoref{sec:lang:watch}                 \\
 \lstinline"wchar_t" & Reserved                                 &
\lstinline"whenever" & \autoref{sec:lang:whenever}              \\
   \lstinline"while" & \lstinline|while&| and \lstinline-while|- flavors &
     \lstinline"xor" & Synonym for \lstinline|^|                \\
  \lstinline"xor_eq" & Synonym \lstinline|^=|                   &
\\
  \hline
\end{tabular}
\else
\begin{tabular}{|c|l|}
  \hline
  Keyword & Remark \\
  \hline
     \lstinline"and" & Synonym for \verb|&&|                    \\
  \lstinline"and_eq" & Synonym for \verb|&=|                    \\
     \lstinline"asm" & Reserved                                 \\
  \lstinline"assert" & \autoref{sec:lang:assert}                \\
      \lstinline"at" & \autoref{sec:lang:at}                    \\
    \lstinline"auto" & Reserved                                 \\
  \lstinline"bitand" & \autoref{sec:lang:op:bitwise}            \\
   \lstinline"bitor" & \autoref{sec:lang:op:bitwise}            \\
    \lstinline"bool" & Reserved                                 \\
   \lstinline"break" & \autoref{sec:lang:break}                 \\
    \lstinline"call" & \autoref{sec:lang:rv}                    \\
    \lstinline"case" & \autoref{sec:lang:switch}                \\
   \lstinline"catch" & \autoref{sec:lang:catch}, \autoref{sec:lang:timeout} \\
    \lstinline"char" & Reserved                                 \\
   \lstinline"class" & \autoref{sec:lang:class}                 \\
 \lstinline"closure" & \autoref{sec:lang:closure}               \\
   \lstinline"compl" & \autoref{sec:lang:op:bitwise}            \\
   \lstinline"const" & \autoref{sec:lang:prop:std}              \\
\lstinline"const_cast" & Reserved                                 \\
\lstinline"continue" & \autoref{sec:lang:continue}              \\
 \lstinline"default" & \autoref{sec:lang:switch}                \\
  \lstinline"delete" & Reserved                                 \\
  \lstinline"detach" &                                          \\
  \lstinline"disown" &                                          \\
      \lstinline"do" & \autoref{sec:lang:do}                    \\
  \lstinline"double" & Reserved                                 \\
\lstinline"dynamic_cast" & Reserved                                 \\
    \lstinline"else" & \autoref{sec:lang:if}                    \\
    \lstinline"enum" & \autoref{sec:lang:continue}              \\
   \lstinline"every" & \autoref{sec:lang:every}                 \\
\lstinline"explicit" & Reserved                                 \\
  \lstinline"export" & Reserved                                 \\
  \lstinline"extern" & Reserved                                 \\
\lstinline"external" & Reserved                                 \\
   \lstinline"float" & Reserved                                 \\
     \lstinline"for" & \verb|for&| and \verb-for|- flavors      \\
 \lstinline"foreach" & Deprecated, use \verb|for|               \\
\lstinline"freezeif" &                                          \\
  \lstinline"friend" & Reserved                                 \\
\lstinline"function" & \autoref{sec:lang:function}, \autoref{sec:lang:closure} \\
    \lstinline"goto" & Reserved                                 \\
      \lstinline"if" & \autoref{sec:lang:if}                    \\
      \lstinline"in" & \autoref{sec:lang:foreach}               \\
  \lstinline"inline" & Reserved                                 \\
     \lstinline"int" & Reserved                                 \\
\lstinline"internal" & Deprecated                               \\
    \lstinline"long" & Reserved                                 \\
    \lstinline"loop" & \verb|loop&|, \verb-loop|-, \autoref{sec:lang:loop} \\
 \lstinline"mutable" & Reserved                                 \\
\lstinline"namespace" & Reserved                                 \\
     \lstinline"new" & \autoref{sec:tut:ctor}                   \\
     \lstinline"not" & Synonym for \verb|!|                     \\
  \lstinline"not_eq" & Synonym for \verb|!=|                    \\
 \lstinline"onleave" & \autoref{sec:lang:at}                    \\
      \lstinline"or" & Synonym for \verb-||-                    \\
   \lstinline"or_eq" & Synonym for \verb-|=-                    \\
 \lstinline"private" & Ignored                                  \\
\lstinline"protected" & Ignored                                  \\
  \lstinline"public" & Ignored                                  \\
\lstinline"register" & Reserved                                 \\
\lstinline"reinterpret_cast" & Reserved                                 \\
  \lstinline"return" & \autoref{sec:lang:rv}                    \\
   \lstinline"short" & Reserved                                 \\
  \lstinline"signed" & Reserved                                 \\
  \lstinline"sizeof" & Reserved                                 \\
  \lstinline"static" & Deprecated                               \\
\lstinline"static_cast" & Reserved                                 \\
  \lstinline"stopif" &                                          \\
  \lstinline"struct" & Reserved                                 \\
  \lstinline"switch" & \autoref{sec:lang:switch}                \\
\lstinline"template" & Reserved                                 \\
    \lstinline"this" &                                          \\
   \lstinline"throw" & \autoref{sec:lang:throw}                 \\
 \lstinline"timeout" & \autoref{sec:lang:timeout}               \\
     \lstinline"try" & \autoref{sec:lang:catch}                 \\
 \lstinline"typedef" & Reserved                                 \\
  \lstinline"typeid" & Reserved                                 \\
\lstinline"typename" & Reserved                                 \\
   \lstinline"union" & Reserved                                 \\
\lstinline"unsigned" & Reserved                                 \\
   \lstinline"using" & Reserved                                 \\
     \lstinline"var" & \autoref{sec:lang:var}, \autoref{sec:lang:slots} \\
 \lstinline"virtual" & Reserved                                 \\
\lstinline"volatile" & Reserved                                 \\
\lstinline"waituntil" & \autoref{sec:lang:waituntil}             \\
   \lstinline"watch" & \autoref{sec:lang:watch}                 \\
 \lstinline"wchar_t" & Reserved                                 \\
\lstinline"whenever" & \autoref{sec:lang:whenever}              \\
   \lstinline"while" & \verb|while&| and \verb-while|- flavors  \\
     \lstinline"xor" & Synonym for \verb|^|                     \\
  \lstinline"xor_eq" & Synonym \verb|^=|                        \\
  \hline
\end{tabular}
\fi

  \caption{Keywords}
  \label{tab:keywords}
\end{table}
\renewcommand{\baselinestretch}{1}



\let\sectionOrig\section
\renewcommand{\section}[1]{\clearpage\sectionObject{#1}}
\section{Pair}

A \dfn{pair} is a container storing two objects, similar in spirit to
\lstinline|std::pair| in \Cxx.


\begin{itemize}
\item \lstinline|asString|\\
  Generate the string \samp{(\var{first}, \var{second})} using
  \code{asPrintable} to convert members to strings.

\item \lstinline|first|\\
  Return the first member of the pair.

\item \lstinline|init(\var{first}, \var{second})|~\\
  Instantiate a new \lstinline|Pair| containing \var{first} and
  \var{second}.

\item \lstinline|second|\\
  Return the second member of the pair.

\item \lstinline|'[]'(\var{index})|\\
  Return the \var{index}-th element.  \var{index} must be 0 or 1.

\item \lstinline|'[]='(\var{index}, \var{value})|\\
  Set (and return) the \var{index}-th element to \var{value}.
  \var{index} must be 0 or 1.

\item \lstinline|'<'(\var{other})|\\
  Lexicographic comparison between two pairs.

\item \lstinline|'=='(\var{other})|\\
  Whether \lstinline|this| and \lstinline|other| have the same
  contents (equality-wise).
\end{itemize}

See below for an example.

\begin{urbiscript}
var p = Pair.new(1, 2);
[00000001] (1, 2)
p.first - p.second;
[00000002] -1
p[0] = 10 | p;
[00000002] (10, 2)
assert(Pair.new(0, 0) < Pair.new(0, 1));
assert(Pair.new(0, 0) < Pair.new(1, 0));
assert(Pair.new(0, 1) < Pair.new(1, 0));
assert(Pair.new(1, 2) == Pair.new(1, 2));
assert(!(Pair.new(1, 1) == Pair.new(2, 2)));
\end{urbiscript}


%%% Local Variables:
%%% mode: latex
%%% TeX-master: "../urbi-sdk"
%%% End:

\let\section\sectionOrig


\part{Second Part}
\label{sec:second}
This is the second part.

\chapter{Chapter II.1}

You want to use \refSlot[Foo]{foo}, and possibly \refSlot[Bar]{foo}.
\settocbibname{\label{sec:bib}The Bibliography}
\setindexname{\label{sec:index}The Index}
\bibliographystyle{plain}
\bibliography{comp.lang.urbi}
\nocite{*}

\chapterIndex
\end{document}

%%% Local Variables:
%%% mode: latex
%%% coding: utf-8
%%% TeX-master: t
%%% ispell-dictionary: "american"
%%% ispell-personal-dictionary: "urbi.dict"
%%% fill-column: 76
%%% End:
